\section{The rational polar completion and the boundary-uniform reciprocal necessity}
\label{sec:polar-cluster-obstruction}

\paragraph{Status.}
The complete author and independent root whole-proof readings
have passed; their exact check scopes are recorded separately.
The two actual polar source coordinates are retained throughout;
the statement does not replace the common source by a bare
diagonal component.

Keep $S=\{\infty\}$, the real smooth normalized step $\phi$,
and the actual Gaussian sections of MA-MO-016.  In its coordinates
write
\[
\begin{gathered}
 H=H_{\Sigma,a,s},\qquad
 \mathcal W=H\oplus\mathbb C^2,\qquad 0<a<1/2,\ s\ge0,\\
 \mathfrak s(u,c)=i_-u+c_0s_0+c_1s_1,\qquad
 s_j=i_+f_j,\qquad i_-u=(u,0)\text{ in these coordinates},\\
 D=P_+-P_-,\qquad
 T_g^0=\mathscr D_g^{\rm raw}i_-=C_gD i_-:
                                  \mathcal E\longrightarrow\mathcal W,\\
 A(w)=\zeta(1/2+w),\qquad p=(s+1)/2,\qquad
 0<\delta_0<\min\{1,(1/2-a)/2\},\qquad
 \mathcal D=C_c^\infty([1,2]).
\end{gathered}
\tag{MA-MO.1300}
\]
The norm of $H$ is exactly the inherited sum norm with weight
$(e^{2ax}+e^{-2ax})(1+x^2)^{s/2}$, and $i_-$ is isometric.
The smooth map $T_g^0$ is the original negative source column,
before any bounded extension is assumed.

Assume for the conditional comparison that all nontrivial zeta
zeros satisfy $|\Re\rho-1/2|\le a$.  Define
\[
 \begin{gathered}
 (\mathrm{BR}_{a,s,\zeta}):\quad
 \exists C,N<\infty\quad
 \delta^p|\zeta(1/2+a+\delta+iT)|^{-1}
                    \le C(2+|T|)^N\\
 \hfill(T\in\mathbb R,\quad0<\delta\le\delta_0).
 \end{gathered}
 \tag{MA-MO.1301}
\]

\paragraph{MA-MO-034 (the actual polar negative column requires reciprocal control).}
If every complex $g\in\mathcal D$ gives a bounded extension
$T_g:H\to\mathcal W$, then (MA-MO.1301) holds.  The same is
true if boundedness is assumed only for every real test in
$\mathcal D$.  In particular the conclusion is necessary for
boundedness of the full original matrix on $\mathcal W$ for
that test family.

If (MA-MO.1301) fails, there is a fixed sequence of compact smooth
inputs $e_n$ of $H$-norm one such that
\[
 \left\{g\in\mathcal D:
       \sup_n\|\mathscr D_g^{\rm raw}i_-e_n\|_{\mathcal W}=\infty\right\}
       \quad\text{is a dense }G_\delta.
 \tag{MA-MO.1302}
\]
This holds separately for the real and complex test spaces.
Each input $i_-e_n$ and every displayed output lies in the
original smooth common source.  Both polar output coordinates
are retained in its norm.

The explicitly hypothetical cluster condition
\[
 \begin{gathered}
 |T_n|\to\infty,\quad d_n\downarrow0,\quad
 0<d_n<\min\{\delta_0,4a/3\},\\
 \#\left\{\rho:\left|\rho-1/2-(a-d_n+iT_n)\right|
                           \le d_n/2\right\}\ge k_n,\\
 \frac{(k_n-p)\log(1/d_n)}{\log(2+|T_n|)}\longrightarrow\infty
 \end{gathered}
 \tag{MA-MO.1303}
\]
implies failure of (MA-MO.1301); the count includes multiplicities.
Every clustered zero is strictly interior.  No such zero
configuration, or independent arithmetic failure of the bound,
is asserted.  The zeta functional equation transfers a cluster
near the left edge to a right-edge cluster of the same form.

\begin{proof}
\emph{1. The actual completion, source and output functional.}
For a compact smooth $v$ use the bounded functional
\[
 \Lambda_v(u,c)=\int v(x)
       \{u(x)+c_0\Theta f_0(x)+c_1\Theta f_1(x)\}\,dx.
 \tag{MA-MO.1304}
\]
It integrates the first actual pair component of $\mathfrak s$.
Weighted Cauchy--Schwarz handles $u$; the fixed section functions
are smooth on the compact support of $v$.
In particular, if $v_T=e^{iTx}v_0$ for fixed compact $v_0$, then
\[
 \sup_T\|\Lambda_{v_T}\|_{\mathcal W^*}<\infty.
 \tag{MA-MO.1305}
\]
For the diagonal part the norm is independent of $T$; for each
polar part the absolute value of its compact integral is bounded
by $\int|v_0\Theta f_j|$.  No polar exponential is put in $H$.

For a compact smooth $y$, the original inverse and its actual
additive source are
\[
 \begin{aligned}
 F_y(x)&=\sum_{n\ge1}\mu(n)n^{-1/2}y(x+\log n),\\
 u_y(t)&=\frac{q(\log|t|)F_y(\log|t|)}{2\sqrt{|t|}},
          \qquad q=1-\phi,\\
 h_y&=u_y+\mathfrak F u_{Ry}.
 \end{aligned}
 \tag{MA-MO.1306}
\]
The inverse sum is locally finite and has a finite right support
endpoint.  The cutoff kills its left end, so $u_y$ is an actual
compact annular even Schwartz function.  Thus $i_+h_y-i_-y$
belongs to the original common-sum space.  Its coordinates in
$\mathcal W$ are the unique actual source-section coordinates
of MA-MO-016, not a discarded-pole representative.
The exact radial normalization is
\[
 \Theta u_y(x)=\sum_{n\ge1}n^{-1/2}
                     q(x+\log n)F_y(x+\log n).
 \tag{MA-MO.1307}
\]
For $y\in\mathcal E$ the same construction is valid using
the rapid right tail, as established in MA-MO-026.

Fix nonzero nonnegative compact smooth functions $h_0,g_0,v_0$ as in
(MA-MO.1204): $g_0$ is supported in $(1,2)$, and
$k_0=C_{g_0^\vee}v_0$ is nonzero, nonnegative, and supported
in an interval $J\subset(r-\log2,r)$ on which $\phi$ has
one nonzero sign, where $r=\sup\operatorname{supp}\phi$.
Set $h_T=e^{-iTx}h_0$, $g_T=e^{-iTx}g_0$, $v_T=e^{iTx}v_0$.
The source dictionary is then
\[
 \begin{gathered}
 v_{g_T}=C_{g_T^\vee}v_T=e^{iTx}k_0(x),\\
 W_T(y)=\sum_{n\ge1}n^{-1/2}v_{g_T}(y-\log n),\qquad
 H_T(w)=\int e^{wx}h_T(x)\,dx,\\
 V_T(w)=\int e^{wx}v_{g_T}(x)\,dx,
 \qquad B_T(w)=\int\phi(y)W_T(y)e^{-wy}\,dy.
 \end{gathered}
 \tag{MA-MO.1308}
\]
The transpose is bilinear and introduces no conjugate of $g_T$.
Every nonunit integer shifts $J$ past $r$, so only $n=1$
contributes to $B_T$.  At $w=a+\delta+iT$,
\[
 \begin{aligned}
 H_T(w)&=\int h_0(x)e^{(a+\delta)x}\,dx,\\
 B_T(w)&=\int_J\phi(x)k_0(x)e^{-(a+\delta)x}\,dx,\\
 V_T(-w)&=\int_Jk_0(x)e^{-(a+\delta)x}\,dx.
 \end{aligned}
 \tag{MA-MO.1309}
\]
The first two have absolute values bounded below uniformly
for $0<\delta\le\delta_0$, and all three are bounded above.

Choose $t_0>0$ by the fixed support conditions (MA-MO.561).
Then $qF_{RU_th_T}=0$ for $t\ge t_0$, before applying
Fourier transformation.  Also the subtracted input has zero
pairing with $v_{g_T}$.  Equations (MA-MO.1306)--(MA-MO.1308)
therefore give, on the actual source,
\[
 \begin{aligned}
 f_T(t)&=\Lambda_{v_T}(T_{g_T}^0U_th_T)\\
       &=\int q(y)W_T(y)F_{h_T}(y-t)\,dy
                                      \quad(t\ge t_0).
 \end{aligned}
 \tag{MA-MO.1310}
\]
This evaluates the represented first component and retains its
two source polar coordinates through $\Lambda_{v_T}$.

\emph{2. The actual orbit limit and the pole-cleared identity.}
The literal wide Euler calculation of (MA-MO.658) yields
\[
 \mathfrak F_T(w)=H_T(w)V_T(-w)
                  -\frac{H_T(w)B_T(w)}{A(w)}-\mathcal R_T(w),
                    \qquad \Re w>1/2,
 \tag{MA-MO.1311}
\]
where $\mathfrak F_T$ is the Laplace integral of (MA-MO.1310)
from $t_0$ to infinity.  The remainder is the compact double
integral (MA-MO.564)--(MA-MO.565) with positive integers.
Its domain is independent of the modulation, and its finite sums
have uniformly bounded moduli.  Therefore
\[
       |\mathcal R_T(a+\delta+iT)|\le C
                         \qquad(T\in\mathbb R,\ 0<\delta\le\delta_0).
 \tag{MA-MO.1312}
\]
It is entire in $w$.  Clearing the zeta pole with
$\widetilde A(w)=(w-1/2)A(w)$ gives the pole-cleared identity
on the wide half-plane
\[
 \widetilde A(w)[H_T(w)V_T(-w)-\mathcal R_T(w)-\mathfrak F_T(w)]
                       =(w-1/2)H_T(w)B_T(w).
 \tag{MA-MO.1313}
\]
The pole of $A$ supplies a removable zero of its reciprocal.

For $w=a+\delta+iT$ define the finite original negative inputs
\[
 \begin{aligned}
 u_{\delta,T,L}&=\int_{t_0}^L e^{-wt}U_th_T\,dt,\\
 u_{\delta,T,L}(x)
      &=e^{-wx}\int_{x-L}^{x-t_0}e^{(a+\delta)y}h_0(y)\,dy.
 \end{aligned}
 \tag{MA-MO.1314}
\]
They are compact and smooth.  The second formula shows that
their moduli are bounded by the infinite-orbit modulus; that
vector has a uniformly bounded compact middle part and exact
right tail $H_T(w)e^{-wx}$.  After increasing $t_0$ so that
the tail starts beyond one, the actual sum weight gives
\[
 \|i_-u_{\delta,T,L}\|_{\mathcal W}
       =\|u_{\delta,T,L}\|_H\le C\delta^{-p}.
 \tag{MA-MO.1315}
\]
Indeed its tail square integral is bounded by a constant times
$\int_1^\infty e^{-2\delta x}(1+e^{-4ax})(1+x^2)^{s/2}\,dx$.
The bound is independent of $T,L$, and uses the explicit tail,
not the weaker triangle estimate on the orbit integral.

For each fixed $\delta$, put $b=a+\delta/2<1/2$.
The closed-strip premise implies the strict margin required by
MA-MO-018 on $\mathcal W_{b,s}$, with $a_*=a$.
The finite orbit inputs converge there, and the actual source
matrix has its proved bounded extension there.  The functionals
$\Lambda_{v_T}$ represent the same compact first-component
integral in every such completion.  Consequently their actual
scalar outputs converge.  The same fixed-width bounds give
holomorphy of the paired Laplace integral on $\Re w>a$.
The identity (MA-MO.1313) continues there from the wide half-plane;
equivalently (MA-MO.1311) continues with the reciprocal's
removable value at $w=1/2$.
No uniform estimate as $b\downarrow a$ is being assumed.
At the displayed $w$, it follows that
\[
 \begin{gathered}
 \Lambda_{v_T}(T_{g_T}^0u_{\delta,T,L})\longrightarrow
                  \mathfrak F_T(a+\delta+iT),\\
 |\mathfrak F_T(a+\delta+iT)|
                  \ge c|A(a+\delta+iT)|^{-1}-C.
 \end{gathered}
 \tag{MA-MO.1316}
\]
Every finite input/output in this limit is a genuine smooth
common-source vector.  Its nonpolar and polar coordinates are
not estimated separately or silently canceled from the norm.

\emph{3. The fixed-support family and Baire's theorem.}
For a fixed compact smooth $u$, the dependence
$g\mapsto T_g^0u$ is continuous into $\mathcal W$ without
assuming a bounded matrix on arbitrary inputs.
Indeed $D i_-u$ is one fixed actual smooth common-source vector,
and $C_g$ is its integrated representation.
The exact representation (MA-MO.259)--(MA-MO.261) is strongly
continuous and locally bounded on $\mathcal W$, including
its theta cocycle and both polar characters.  Thus on a fixed
compact support its integral has norm bounded by a constant
times $\|g\|_{L^1}$ times that vector's norm.

If every complex test has a bounded column, the closed sets
\[
 \mathcal F_m=\{g\in\mathcal D:
       \|T_g^0u\|_{\mathcal W}\le m\|u\|_H
                            \text{ for every }u\in C_c^\infty\}
 \tag{MA-MO.1317}
\]
cover $\mathcal D$.  Baire's theorem, subtraction and scaling
give finite $N,C$ with
$\|T_g\|_{H\to\mathcal W}\le C\|g\|_{C^N}$.
If the hypothesis concerns only real tests, first apply Baire
over the real Fr\'echet space and then complexify this bound.
Since $\|g_T\|_{C^N}\le C_N(2+|T|)^N$, equations
(MA-MO.1305), (MA-MO.1315) and (MA-MO.1316) give
\[
 c|A(a+\delta+iT)|^{-1}-C
          \le C'(2+|T|)^N\delta^{-p}.
 \tag{MA-MO.1318}
\]
This proves the required estimate (MA-MO.1301).

\emph{4. The fixed genuine-input sequence under failure.}
If (MA-MO.1301) fails, choose pairs $(T_n,\delta_n)$ for which
$\delta_n^p|A(a+\delta_n+iT_n)|^{-1}$ exceeds an arbitrarily
large fixed multiple of $n(2+|T_n|)^n$.
For each pair choose a finite $L_n$ so that the scalar error
in (MA-MO.1316) is at most one.  Normalize
$e_n=u_{\delta_n,T_n,L_n}/\|u_{\delta_n,T_n,L_n}\|_H$.
The denominator is nonzero, and the preceding bounds imply
\[
           \|T_{g_{T_n}}^0e_n\|_{\mathcal W}
                               \ge n(2+|T_n|)^n.
 \tag{MA-MO.1319}
\]
The input has zero source polar coordinates; its output retains
both coordinates.  The fixed sequence is chosen through finite
actual orbit integrals, not through a proposed completed inverse.

For these fixed $e_n$, each set
$\{g:\sup_n\|T_g^0e_n\|_{\mathcal W}\le m\}$ is closed.
If it had interior, the same subtraction/scaling argument would
give a finite $C^N$ seminorm bound, contradicting (MA-MO.1319).
For the real space the bound complexifies before evaluating the
complex modulations.  These sets are nowhere dense, so their
complement is the asserted dense $G_\delta$ in either space.
Unboundedness of this column excludes a bounded full matrix.

\emph{5. The conditional zeta cluster.}
At large height the finite disk centered at $2+iT_n$ with radius
between $12$ and $13$ contains no zeta pole.  Divide zeta by its
finite disk Blaschke product.  The polynomial strip upper bound
and maximum principle bound the remaining holomorphic factor
by $C(2+|T_n|)^D$.  At the raw evaluation point
$1/2+a+d_n+iT_n$, every clustered factor has modulus at most
$C_1d_n$, and every other factor has modulus at most one.
Thus, exactly as in (MA-MO.1219)--(MA-MO.1221),
\[
 d_n^p|\zeta(1/2+a+d_n+iT_n)|^{-1}
       \ge C^{-1}(2+|T_n|)^{-D'}d_n^{p-k_n}.
 \tag{MA-MO.1320}
\]
The ordinary fixed-window count $k_n=O(\log(2+|T_n|))$
absorbs the constant $C_1^{k_n}$ in a power of the height.
Condition (MA-MO.1303) now forces failure of (MA-MO.1301).
All source conclusions were proved before making this conditional
specialization.  No existence of such a cluster is inferred.
\end{proof}

\paragraph{Exact relation to the arithmetic and operator gates.}
The theorem concerns the actual negative column into the full
polar completion, and so supplies a necessary condition for the
original full matrix.  It is not a statement about the bounded
positive source return, the quotient-return norms, or an
orthogonal-PW identification.  The pole at one is retained through
the actual source sections and the removable inverse value; it
does not exempt the full matrix from the reciprocal estimate.
Neither this estimate nor full individual jet retention is
asserted sufficient at arbitrary widths.

The user-stipulated A supplies zeta RH.  Then MA-MO-018 gives
the positive fixed-gap branch for every $0<a<1/2$, $s\ge0$
and every compact smooth test, and (MA-MO.1301) follows.
This explicitly named zero-location input is not replaced by
an independent arithmetic reciprocal estimate in this theorem.
No new arithmetic movement of FT002 is claimed.

\paragraph{Source and finite-length placement.}
The actual polar-section completion, local representation and
negative embedding are (MA-MO.250)--(MA-MO.261).  The source
dictionary, bounded compact output functional and pole-cleared
Laplace identity are (MA-MO.653)--(MA-MO.659), proved from
Meyer's original common-space maps, supplied \texttt{main.tex}
lines 3808--4043.  The wider-width source agreement is MA-MO-018,
and the exact norm and cluster analysis have been written above.
No new full primary-source proof reading is claimed here.
In the inherited cote~1 framework, BR-TR-01/02, scans 079--080
and 111--112, the two actual polar sections and finite-length
divisor modules remain separate data.  Their $K_0$ classes and
finite-dimensional exterior-power determinants do not supply the
additional uniform source norm required for a nuclear trace.
